\cotefile{14}
\scanpage{45}
\mspage{2.}
$i\handwritten{\geqslant} n$, surjectif pour $i=n-1$. En fait, sauf erreur je sais prouver la conj.$A$ à partir de la forme plus faible suivante (dans la nature d'un théorème de Lefschetz ``faible'') :

\uline{Conjecture} $A'$. Pour tout entier $i$ tel que $2i\handwritten{\geqslant} n-1$, l'application image directe $C^i(Y)\handwritten{\longrightarrow} C^{i+1}(X)$ est surjective, i.e. tout cycle algébrique de dimension $j\handwritten{<}n/2$ sur $X$ est $\handwritten{\tau}$-équivalent à l'image d'un cycle algébrique à coeff. rationnels de $Y$.

Ici encore il suffit sauf erreur de faire $n=2i+1$. D'ailleurs, dans l'argument utilisé pour ramener $A$ à $A'$, il semble qu'il faille démontrer $A'$ sur un corps de base pas nécessairement algébriquement clos, et pour des classes de cycles rationnelles sur ce dernier, \uline{ou} se borner à énoncer $A$ pour l'équivalence cohomologique (mais alors $A$ perd le caractère ``purement géométrique'' que je t'avais promis !).

Comme je te l'avais dit, $A$ implique ``la formule de Künneth pour les cycles'' (j'avais d'ailleurs charrié en prétendant que les deux énoncés sont équivalents). Par là, il implique tous les théorèmes d'entiers qu'on peut désirer (par exemple pour des coefficients de polynômes caractéristiques), à cela près qu'il semble possible qu'il y ait des puissances de $p=\operatorname{car} k$ en dénominateur. Cela implique donc les conjectures de Weil, à la seule exception de la question des valeurs absolues des valeurs propres, qui sera justiciable de $B$. Je te signale aussi que c'est la conjecture $A$ qui semble le minimum minimorum pour donner une définition en forme utilisable de la notion de motif sur un corps.

\uline{Conjecture} $B$\quad Supposons $n=2m$, et soit $P^m(X)$ le noyau de l'homomorphisme de multiplication par $\handwritten{\xi}$ de $C^m(X)$ dans $C^{m+1}(X)$. Alors la forme $(-1)^m\handwritten{\unsure{\varepsilon|\mathcal{E}}}(xx')$ sur $P^m(X)$ est définie positive.

Pour les conjectures de Weil, il suffirait d'ailleurs de prouver que cette forme est positive. Mais la forme forte implique aussi d'autres résultats agréables, comme le fait que $\handwritten{\tau}$-équivalence $=$ équivalence numérique ($=$ équivalence homologique à coeff $\underline{\mathrm{Q}}_{\unsure{l|1|\ell}}$, cette dernière étant coincée entre les deux) et le fait que les $C^m(X)$ sont de dimension finie,
% archive sheet number: 44
